\usepackage[utf8]{inputenc}
\usepackage[T1]{fontenc}
\usepackage[sfdefault]{carlito}



\usepackage{animate}
\usepackage{tikz}
\usetikzlibrary{positioning,chains,fit,shapes,calc}
\usetikzlibrary{shapes,arrows}
%\usepackage{amsmath,bm,times}
\newcommand{\mx}[1]{\mathbf{{#1}}} % Matrix command
\newcommand{\vc}[1]{\mathbf{{#1}}} % Vector command


\usepackage{bbm, dsfont}
%\RequirePackage{mathspec} 

\usetikzlibrary{decorations.pathmorphing} % noisy shapes
\usetikzlibrary{fit}					% fitting shapes to coordinates
\usetikzlibrary{backgrounds}	% drawing the background after the foreground
\usetikzlibrary{shadows}
\usetikzlibrary{mindmap}
%%

\usepackage{tabularx}   % Erweiterte Tabellen Optionen
\usepackage{booktabs}
\usepackage{multicol}
\usepackage{multirow}

\usepackage{xcolor}
\usepackage{scalefnt}
\usepackage{amsmath,amsfonts,dsfont}

\usepackage{tikz}

\usepackage{color}

%\usepackage{tcolorbox}% http://ctan.org/pkg/tcolorbox

\definecolor{arrowblue}{RGB}{98,145,224}
\definecolor{yellowvirtual}{RGB}{248,216,90}
\definecolor{myblueauthor}{rgb}{0.215,0.625,0.873}


\definecolor{mycolor}{rgb}{0.122, 0.435, 0.698}% Rule colour
\colorlet{pink}{black!40}
\colorlet{blue}{cyan!60}

\colorlet{insidecode}{pink!30}
\colorlet{myframeinside}{pink!30}

\colorlet{mybrown}{cyan!30}
\colorlet{myframe}{cyan!30}

\newcounter{myexamplectr}
 \makeatletter
  \@addtoreset{myexamplectr}{section}
  \makeatother
  
%\definecolor{mybluetitle}{rgb}{0.0, 0.2, 0.45}
 
\definecolor{myblueexercise}{rgb}{0.215,0.625,0.873}
  
%\newenvironment{sjexample}[2]{%
% \refstepcounter{myexamplectr}% increment the environment's counter
%% \begin{block}{}\centering
%%\textsc{Exemplo \themyexamplectr}
%%\end{block}
%\hspace*{-0.4cm}\begin{tikzpicture}[scale=1.0, auto]
%%\useasboundingbox (-1,-1) rectangle (5,4);
%\coordinate (a) at (0,0);
%%\coordinate (a1) at (0,-0.3);
%\begin{scope}
%\draw[decoration=penciline, decorate,ultra thick,yellow!70!black,rotate=180,xshift=1.8cm,yshift=0.4cm] (0,0) arc (0:-90:1cm);
%				\node[noise,rectangle,minimum height=0.6cm] (a) at (a) {\begin{minipage}{3.0cm}\large \centering Exemplo~\themyexamplectr \end{minipage}};
%\end{scope}
%\end{tikzpicture}
%
%%\emRed{Exemplo \themyexamplectr}
%%\numberwithin{myexamplectr}{section}
%#1
%\def\temp{#2}\ifx\temp\empty
%%
%\else
%  \emGray{Solução}%
%\fi
%
%\vspace*{-1.0cm}
%\hspace*{0.95\linewidth}\begin{tikzpicture}[decoration=penciline, decorate]%,remember picture,overlay]
%\begin{scope}
%\draw[decorate,ultra thick,yellow!70!black] (0,0) arc (0:-90:1cm);
%\end{scope}
%%\hfill${\triangleleft}$
%\end{tikzpicture}
%}


\usepackage[most]{tcolorbox}

\newtcbtheorem[auto counter]{sjexample}{\bfseries Exemplo}{enhanced,%,drop shadow={black!50!white},
% width = 1.02\linewidth,
  coltitle=white,
  top=0.2in,
  coltext=myblue,
%  colback=green!5,
  colframe=myblueexercise!70,
  fonttitle=\bfseries,
  attach boxed title to top left=
  {xshift=10pt,yshift=-\tcboxedtitleheight/2},
  boxed title style={sharp corners, 
                     size=small,colback=myblueexercise!70, 
                     borderline={0.15mm}{-0.15mm}{myblueexercise!70,solid},
                     coltext=black},
  colback=myblueexercise!5
}{sjexample}

\newtcbtheorem[auto counter]{sjsolution}{\bfseries \textsc{Solução}}{enhanced,%,drop shadow={black!50!white},
 width = 1.0\linewidth,
  coltitle=white,
  top=0.2in,
  coltext=myblue,
%  colback=green!5,
  colframe=yellowvirtual!90!black,
%  fonttitle=\bfseries\sc,
  left skip={.2\linewidth},
  right skip={.2\linewidth},
  attach boxed title to top center=
  {xshift=0pt,yshift=-\tcboxedtitleheight/2},
  boxed title style={sharp corners, 
                     size=small,colback=yellowvirtual!90!black, 
                     borderline={0.15mm}{-0.15mm}{yellowvirtual!90!black,solid},
                     coltext=black},
  colback=yellowvirtual!15
}{sjsolution}



\resetcounteronoverlays{tcb@cnt@sjexample}
\newcommand{\solucao}{\emGray{Solução}~~~~}
\newcommand{\prova}{\emGray{Prova}~~~~}

% Define box and box title style
\tikzstyle{mybox} = [draw=myblue, fill=white, very thick,
    rectangle, rounded corners, inner sep=10pt, inner ysep=10pt]
\tikzstyle{fancytitle} =[fill=red, text=white]




\newenvironment{sgexample}[2]{%
% \refstepcounter{myexamplectr}% increment the environment's counter
% \stepcounter{myexamplectr}
\vspace*{0.3cm}
\begin{tikzpicture}
\node [mybox] (box){%
    \begin{minipage}{0.95\linewidth} \vspace*{0.2cm}
#1

%\addtocounter{myexamplectr}{0}
\def\temp{#2}\ifx\temp\empty
%
\else
  \emGray{Solução}%
  #2
\fi
    \end{minipage}
};
\node[rectangle,fill=white,very thick,draw=myblue,minimum height=0.6cm, right=10pt] at (box.north west) {\begin{minipage}{2.0cm}\centering Exemplo~\themyexamplectr \end{minipage}};
%\node[fancytitle, rounded corners] at (box.east) {$\clubsuit$};
\end{tikzpicture}%
% \refstepcounter{myexamplectr}% increment the environment's counter
%\hspace*{-0.4cm}\begin{tikzpicture}[scale=1.0, auto]
%\coordinate (a) at (0,0);
%%\coordinate (a1) at (0,-0.3);
%\begin{scope}
%\draw[decoration=penciline, decorate,ultra thick,yellow!70!black,rotate=180,xshift=1.8cm,yshift=0.4cm] (0,0) arc (0:-90:1cm);
%				\node[noise,rectangle,minimum height=0.6cm] (a) at (a) {\begin{minipage}{3.0cm}\large \centering Exemplo~\themyexamplectr \end{minipage}};
%\end{scope}
%\end{tikzpicture}
%
%
%#1
%\def\temp{#2}\ifx\temp\empty
%%
%\else
%  \emGray{Solução}%
%\fi
%
%\vspace*{-1.0cm}
%\hspace*{0.95\linewidth}\begin{tikzpicture}[decoration=penciline, decorate]%,remember picture,overlay]
%\begin{scope}
%\draw[decorate,ultra thick,yellow!70!black] (0,0) arc (0:-90:1cm);
%\end{scope}
%\end{tikzpicture}
}

%\begin{tikzpicture}[remember picture,overlay]
%\coordinate (aa) at ($(-0.25,4.35)$);
%\coordinate (bb) at ($(11.75,0.65)$);
%\node[draw=myframe,line width=1pt,fill=mybg,opacity=0.25,rectangle,rounded corners,fit=(aa) (bb)] {};
%%\coordinate (cc) at ($(-4.2,-0.4)$);
%%\coordinate (dd) at ($(-0.15,-0.9)$);
%%\node[text=black,line width=1pt,rectangle,rounded corners,fit=(cc) (dd)] {$K_{4}$ -- 4 faces};
%\end{tikzpicture}

%\newcounter{myalgctr}
%
%\newenvironment{myalg}{%      define a custom environment
%   \bigskip\noindent%         create a vertical offset to previous material
%   \refstepcounter{myalgctr}% increment the environment's counter
%   \textsc{Algorithm \themyalgctr}% or \textbf, \textit, ...
%   \newline%
%   }{\par\bigskip}  %          create a vertical offset to following material
%\numberwithin{myalgctr}{section}

\AtBeginEnvironment{sgexample}{\refstepcounter{myexamplectr}}


%\newcommand\myframe[2]{
%\begin{frame}
%  \frametitle{#1}
%  #2
%\end{frame}}

%\definecolor{BlueLUH}{cmyk}{1.0,0.7,0,0}
\definecolor{mybg1}{RGB}{0,115,207}
\colorlet{mybg}{mybg1!20!white}




%\makeatletter
%\newcommand{\mybox}[1]{%
%  \setbox0=\hbox{#1}%
%  \setlength{\@tempdima}{\dimexpr\wd0+13pt}%
%  \begin{tcolorbox}[colframe=mycolor,boxrule=0.5pt,arc=4pt,
%      left=6pt,right=6pt,top=6pt,bottom=6pt,boxsep=0pt,width=\@tempdima]
%    #1
%  \end{tcolorbox}
%}
%\makeatother

\newcommand\myBox[3]{{\colorbox{#1}{\color{#2} {#3}}}}

\usepackage{ragged2e}

%\usepackage{midlevel-plot}

\usepackage[T1]{fontenc}
\usepackage{mfirstuc}
%\usepackage{fontspec}

\newcommand{\ssc}[1]{\textsc{\capitalisewords{\MakeLowercase{#1}}}}

\definecolor{myblue}{rgb}{0.0, 0.2, 0.45}
\definecolor{myred}{rgb}{0.95, 0.2, 0.0}
\definecolor{mywhite}{rgb}{0.95, 0.95, 0.95}
\definecolor{mygreen}{rgb}{0.0, 0.55, 0.0}
\definecolor{myorange}{rgb}{0.95, 0.5, 0.0}
\definecolor{myyellow}{rgb}{0.85, 0.85, 0.0}
\definecolor{mygray}{rgb}{0.15, 0.15, 0.15}
\definecolor{myblack}{rgb}{0.05, 0.05, 0.05}
\definecolor{myredb}{rgb}{0.75, 0.2, 0.0}
\definecolor{mybggraph}{rgb}{0.52, 0.52, 0.52}


\newcommand\emRed[1]{\myBox{myred!70!myblack}{mywhite}{#1}}
\newcommand\emBlue[1]{\myBox{myblue}{mywhite}{#1}}
\newcommand\emGray[1]{\myBox{mygray}{myblack}{\color{myblack} #1}}
\newcommand\emBlack[1]{\myBox{myblack}{mywhite}{#1}}
\newcommand\emGreen[1]{\myBox{mygreen}{mywhite}{#1}}
\newcommand\emOrange[1]{\myBox{myorange}{mywhite}{#1}}

\newcommand\ie{\em{i.e.}}


\newcommand\parGray[2][0.999\linewidth]{{
\emGray{ 
\begin{minipage}{#1} \centering
#2
\end{minipage}  
  }
  }}
  
  \newcommand\parBlue[2][0.999\linewidth]{{
\emBlue{ 
\begin{minipage}{#1} \centering \rmfamily\scshape
#2
\end{minipage}  
  }
  }}
  
  \newcommand\parRed[2][0.999\linewidth]{{
\emRed{ 
\begin{minipage}{#1} \centering \rmfamily\scshape #2
\end{minipage}  
  }
  }}
    
  \newcommand\parBlack[2][0.999\linewidth]{{
\emBlack{ 
\begin{minipage}{#1} \centering \rmfamily\scshape
#2
\end{minipage}  
  }
  }}
        
    
  \newcommand\parGreen[2][0.999\linewidth]{{
\emGreenDark{
\begin{minipage}{#1} \centering
#2
\end{minipage}  
  }
  }}
  
      \newcommand\parOrange[2][0.999\linewidth]{{
\emOrange{
\begin{minipage}{#1} \centering
#2
\end{minipage}  
  }
  }}
  

%\usetheme{cambridge}


\usetikzlibrary{fadings, backgrounds}
\tikzfading[name=fade out, inner color=transparent!100, outer color=transparent!25]
\tikzfading[name=fade in, inner color=transparent!100, outer color=transparent!0]

\usetikzlibrary{spy}

\usepackage{animate}


\usetikzlibrary{fadings, backgrounds}
\tikzfading[name=fade out, inner color=transparent!100, outer color=transparent!25]
\tikzfading[name=fade in, inner color=transparent!100, outer color=transparent!0]

\usetikzlibrary{spy}


\usepackage{animate}

%\usepackage[utf8]{inputenc}
%\usepackage{fontawesome}


\usetikzlibrary{decorations.text}
\definecolor{mygray}{RGB}{208,208,208}
\definecolor{mymagenta}{RGB}{226,0,116}
\newcommand*{\mytextstyle}{\sffamily\Large\bfseries\color{black!85}}
\newcommand{\arcarrow}[3]{%
   % inner radius, middle radius, outer radius, start angle,
   % end angle, tip protusion angle, options, text
   \pgfmathsetmacro{\rin}{1.7}
   \pgfmathsetmacro{\rmid}{2.2}
   \pgfmathsetmacro{\rout}{2.7}
   \pgfmathsetmacro{\astart}{#1}
   \pgfmathsetmacro{\aend}{#2}
   \pgfmathsetmacro{\atip}{5}
   \fill[mygray, very thick] (\astart+\atip:\rin)
                         arc (\astart+\atip:\aend:\rin)
      -- (\aend-\atip:\rmid)
      -- (\aend:\rout)   arc (\aend:\astart+\atip:\rout)
      -- (\astart:\rmid) -- cycle;
   \path[
      decoration = {
         text along path,
         text = {|\mytextstyle|#3},
         text align = {align = center},
         raise = -1.0ex
      },
      decorate
   ](\astart+\atip:\rmid) arc (\astart+\atip:\aend+\atip:\rmid);
}
\usepackage{animate}

\usepackage{color}
\newcommand{\head}[1]{\textnormal{\textbf{#1}}}
\newcommand{\normal}[1]{\multicolumn{1}{l}{#1}}
%\usepackage{booktabs}
\usepackage{pgfplotstable}
\usepackage{tikz}
%%%%%%%%%%%%%%%%%%%%%%%%%%%%%%%%%%%%%%%%%%%%%%%%%%%%%%%%%%%%%%%%%%%%%%%%%%%%%%%%%%%%%%%%
%% color definitions %%%%%%%%%%%%%%%%%%%%%%%%%%%%%%%%%%%%%%%%%%%%%%%%%%%%%%%%%%%%%%%%%%%
%%%%%%%%%%%%%%%%%%%%%%%%%%%%%%%%%%%%%%%%%%%%%%%%%%%%%%%%%%%%%%%%%%%%%%%%%%%%%%%%%%%%%%%%
%\definecolor{gris1}{RGB}{220,220,220}
%\definecolor{gris2}{RGB}{10,10,10}
%\definecolor{couleurbarretitreslide}{RGB}{128,0,0}
%\definecolor{couleurtitreslide}{RGB}{220,220,220}
%\definecolor{couleurrouge}{RGB}{128,0,0}
%\definecolor{couleurtexteblock}{RGB}{0,0,255}%{5,5,145}
%\definecolor{redB}{RGB}{180,10,10}
%\definecolor{blueB}{RGB}{10,10,180}
%


\newcommand{\RedBold}[1]{{\bf \color{red!70!black} #1}}
\newcommand{\BlueBold}[1]{{\bf \color{blue!40!black} #1}}
\newcommand{\GreenBold}[1]{{\bf \color{green!40!black} #1}}
\newcommand{\BlackBold}[1]{{\bf \color{black} #1}}
\newcommand{\BrownBold}[1]{{\bf \color{brown} #1}}
\newcommand{\CyanBold}[1]{{\bf \color{cyan} #1}}
\newcommand{\MagentaBold}[1]{{\bf \color{magenta} #1}}
\newcommand{\YellowBold}[1]{{\bf \color{yellow} #1}}
\newcommand{\PurpleBold}[1]{{\bf \color{purple} #1}}

\newcommand{\redB}[1]{\RedBold{#1}}
\newcommand{\blueB}[1]{\BlueBold{#1}}
\newcommand{\brownB}[1]{{\BrownBold{#1}}}

\newcommand{\redBZ}[1]{\scalefont{1.2} \redB{#1}}
\newcommand{\blueBZ}[1]{\scalefont{1.2} \blueB{#1}}
\newcommand{\brownBZ}[1]{{\scalefont{1.2}\brownB{#1}}}


\usepackage[absolute,overlay]{textpos}

\newcommand\exactlocation[4]{
 \begin{textblock*}{#1}(#2,#3)
#4
\end{textblock*}
%\null\cleartextposbox 
}

\usepackage{media9}
\usepackage{multimedia}


%
%\definecolor{myblue}{rgb}{0.0, 0.2, 0.45}
%\definecolor{myred}{rgb}{0.95, 0.2, 0.0}
%\definecolor{mywhite}{rgb}{0.95, 0.95, 0.95}
%\definecolor{mygreen}{rgb}{0.0, 0.55, 0.0}
%\definecolor{myorange}{rgb}{0.95, 0.5, 0.0}
%\definecolor{myyellow}{rgb}{0.85, 0.85, 0.0}
%\definecolor{mygray}{rgb}{0.45, 0.45, 0.45}
%\definecolor{myblack}{rgb}{0.05, 0.05, 0.05}
%\definecolor{myredb}{rgb}{0.75, 0.2, 0.0}
%\definecolor{mybggraph}{rgb}{0.52, 0.52, 0.52}

%
%
%\newcommand\parGray[2][0.955\linewidth]{{
%\emGray{ 
%\begin{minipage}{#1} \centering
%#2
%\end{minipage}  
%  }
%  }}
%  
%  \newcommand\parBlue[2][0.955\linewidth]{{
%\emBlue{ 
%\begin{minipage}{#1} \centering
%#2
%\end{minipage}  
%  }
%  }}
%  
%  \newcommand\parRed[2][0.955\linewidth]{{
%\emRed{ 
%\begin{minipage}{#1} \centering
%#2
%\end{minipage}  
%  }
%  }}
%    
%  \newcommand\parBlack[2][0.955\linewidth]{{
%\emBlack{ 
%\begin{minipage}{#1} \centering
%#2
%\end{minipage}  
%  }
%  }}
%        
%    
%  \newcommand\parGreen[2][0.955\linewidth]{{
%\emGreenDark{
%\begin{minipage}{#1} \centering
%#2
%\end{minipage}  
%  }
%  }}
%  
%      \newcommand\parOrange[2][0.955\linewidth]{{
%\emOrange{
%\begin{minipage}{#1} \centering
%#2
%\end{minipage}  
%  }
%  }}


\renewcommand\textsc[1]{{\rmfamily\scshape #1}}

\usepackage{xcolor,colortbl}
%\usepackage{soul}
\usepackage{cancel}

\newcommand{\mc}[2]{\multicolumn{#1}{c}{#2}}
\definecolor{Gray}{gray}{0.85}
\definecolor{LightCyan}{rgb}{0.88,1,1}

\renewcommand<>\cellcolor[1]{\only#2{\beameroriginal\cellcolor{#1}}}

\newcommand\selectedcell[2]{{\cellcolor<#1>{mygray!20!mywhite}\textcolor{myblack}{#2}}}
\newcommand\highlightcell[2]{\cellcolor<#1>{myorange!50!mywhite}\textcolor{myblack}{#2}}
\newcommand\missedcell[2]{\cellcolor<#1>{myblack!50!mywhite}\textcolor{mywhite}{#2}}
\newcommand\ignoredcell[2]{\cellcolor<#1>{myblack!05!mywhite}\textcolor{myblack}{\cancel{#2}}}
\newcommand\unselectedcell[2]{{\cellcolor<#1>{mywhite}\textcolor{myblack}{#2}}}

\usepackage{multirow}
\usepackage{xcolor,colortbl}
\usepackage{hhline}

\newcolumntype{L}[1]{>{\raggedright\let\newline\\\arraybackslash\hspace{0pt}}m{#1}}
\newcolumntype{C}[1]{>{\centering\let\newline\\\arraybackslash\hspace{0pt}}m{#1}}
\newcolumntype{R}[1]{>{\raggedleft\let\newline\\\arraybackslash\hspace{0pt}}m{#1}}


%% MACROS
%\newcommand{\reals}{\ensuremath{\mathbb{R}}}
\newcommand{\integers}{\ensuremath{\mathbb{Z}}}
% macro for a set. extra pair of braces inside \ensuremath to ensure
% the set is typseset in one line.
%\newcommand{\set}[1]{\ensuremath{{\{#1\}}}}
% macro for the cardinality of a set
%\newcommand{\card}[1]{\ensuremath{|#1|}}
%\newcommand{\unionfind}{\textsc{Union-Find}\xspace}

%\newcommand{\calc}{\ensuremath{{\cal C}}\xspace}
%\newcommand{\calo}{\ensuremath{{\cal O}}\xspace}
%\newcommand{\opt}[1]{\ensuremath{\text{OPT}(#1)}}

% flow macros.
\newcommand{\bottleneck}[1]{\ensuremath{\text{bottleneck}(#1)}}
\newcommand{\fin}[1]{\ensuremath{f^{\text{in}}(#1)}}
\newcommand{\fout}[1]{\ensuremath{f^{\text{out}}(#1)}}

% np-completeness macros.
\newcommand{\negation}[1]{\ensuremath{\overline{#1}}}
\newcommand{\reducible}[2]{{\textsc{#1}\ensuremath{~\leq_P~}\textsc{#2}}}
\newcommand{\equivalent}[2]{{\textsc{#1}\ensuremath{~\equiv_P~}\textsc{#2}}}
\newcommand{\calp}{\ensuremath{{\cal P}}\xspace}
\newcommand{\nph}{\ensuremath{{\cal NP}}-Hard\xspace}
\newcommand{\calnp}{\ensuremath{{\cal NP}}\xspace}
\newcommand{\npc}{\calnp-Complete\xspace}
\newcommand{\conp}{\text{co-}\calnp}%\ensuremath{{\cal NP}}\xspace}

\newenvironment{myfont}{\wela\selectfont\color{myred}}{\par}



\tikzset{pencil draw/.style={% code from Claudio Fiandrino's answer at https://tex.stackexchange.com/a/95333/
    draw,
    decorate,
    decoration={random steps, segment length=3mm, amplitude=0.2mm}
  },
}



\tikzset{
>=stealth',
  flimlayer/.style={
	draw=mygreen!60!myblack,     minimum width=0.4em, 
	text = white,
    minimum height=6em, fill=mygreen!60!myblack,line width=1.5pt,rounded corners
    },
  line/.style={draw, thick, <-},
  element/.style={
    tape,
    top color=white,
    bottom color=blue!50!black!60!,
    minimum width=8em,
    draw=blue!40!black!90, very thick,
    text width=10em, 
    minimum height=3.5em, 
    text centered, 
    on chain},
  every join/.style={->, thick,shorten >=1pt},
  decoration={brace},
  tuborg/.style={decorate},
  tubnode/.style={midway, right=2pt},
}







\makeatother


\tikzstyle{noise}=[circle,
                                    thick,
                                    minimum size=1.2cm,
                                    draw=yellow!85!black,
                                    fill=yellow!40,
                                    decorate,
                                    decoration={random steps,
                                                            segment length=2pt,
                                                            amplitude=2pt}]

